% 本文件由 LaTeX 排版 Agent 根据有效 translation.json 自动生成。
% author=Athanasius; work_id=2805; title=论路加福音10:22及玛窦福音11:27

\chapter{论\BibleBook{路加}{福音}10:22及\BibleBook{玛窦}{福音}11:27}

% structure: p0001 source=h1 target=omitted reason=page-title-already-rendered-by-parent

论\BibleBook{路加}{福音}10:22及\BibleBook{玛窦}{福音}11:27\par

§ \Emph{1. 这段经文所指的不是永恒圣言，而是降生成人的那一位。}\par

\BibleQuote{我父将一切交给了我；除了父，没有一个认识子是谁；除了子和子所愿意启示的人，也没有一个认识父是谁。}\par

\Person{亚略}派、\Person{欧瑟比乌斯}及其同党因为不明了这点，而对主肆意不敬。因为他们说，如果一切都被交付了（这里的\Quote{一切}指创造的主权），那么曾经有一段时间祂没有这一切。但如果祂没有这一切，祂就不是出自圣父，因为如果祂是出自圣父，祂就因此常常拥有这一切，而不需要接受了。但这点将更清楚地揭露他们的愚蠢。因为这里所说的，并非指对创造的统治权，也不是指掌管天主的工程，而是为显示降生成人的计划 (\OriginalTerm{经世}{\GreekText{τῆς} \GreekText{οἰκονομίας}}) 的一部分用意。因为如果祂正说话的时候，它们\Quote{被交付}给祂，那么显然在祂接受之前，创造就没有圣言。那么，\BibleQuote{万有都因祂而存在} \BibleRef{哥 1:17} 这句话又如何呢？但如果这一切在创造之初就同时被\Quote{交付}给祂，这样的交付便是多余的，因为\BibleQuote{万有是藉着祂而造成的} \BibleRef{若 1:3}，那些本是主亲手所做的工程，就不需要再交付给祂了。因为在创造它们时，祂就是这些被造之物的主。即便假设它们在被造之后才\Quote{交付}给祂，请看这何等荒谬。因为如果它们\Quote{被交付}，而圣父在交付后便退隐，那么我们就有落入某些人所编造神话的危险：即天主将祂的工程交给了圣子，然后自己离开了。或者，如果圣子拥有它们时圣父也仍然拥有，那么我们应该说，这不是\Quote{被交付}，而是天主接纳子为同工，如同\Person{保禄}接纳\Person{息耳瓦诺}那样。但这更为荒谬；因为天主并非不圆满，祂也没有呼召圣子来帮助自己的所需；而是作为圣言的圣父，祂藉着圣言创造万物，并不把创造交付给祂，而是藉着祂并在祂内对创造施行天主眷顾，以致没有圣父许可，连一只麻雀也不会掉在地上 \BibleRef{玛 10:29}，野草没有天主就无从穿着 \BibleRef{玛 6:30}，而且圣父工作直到如今，圣子也工作直到如今 \BibleRef{若 5:17}。因此，不敬神者的意见是徒劳的。因为经文所指的并非他们所想，而是指降生成人这件事。\par

§2. \Emph{一切交付给降生成人的圣子的意义与目的。}\par

因为人犯了罪，堕落了，且由于他的堕落，万物陷入混乱：死亡从\Person{亚当}直到\Person{梅瑟}一直为王 \BibleRef{罗 5:14}，大地受了诅咒，阴府敞开，乐园关闭，天堂被触怒，人最终败坏且沦为禽兽 \BibleRef{咏 48:12}，魔鬼正对我们耀武扬威——那时，富于慈爱的天主，不愿祂按自己肖像所造的人丧亡，便说：\BibleQuote{我将派遣谁呢？谁肯为我们去呢？} \BibleRef{依 6:8}。众人都默不作声时，圣子说：\Quote{我在这里，请派遣我。}就在那时，圣父说\Quote{去罢}，便将人\Quote{交付}给圣子，好使圣言自己成为血肉，藉着取得血肉，将人性完全恢复。因为人\Quote{被交付}给祂，就像病人托付给医生，为治好毒蛇咬伤；就像死者托付给生命，为使死人复活；就像黑暗托付给光明，为照亮幽暗；又因为祂是圣言，为更新理性本质 (\OriginalTerm{理性本质}{\GreekText{τὸ} \GreekText{λογικόν}})。既然万物都\Quote{被交付}给了祂，而祂成了人，一切便立即改正并趋于完善。大地获得祝福代替诅咒，乐园向右盗敞开，阴府战栗，坟墓敞开死者复活，天门为迎接那\Quote{来自厄东}的而高耸 \BibleRef{咏 24:7}。看，救主自己清楚说明\Quote{万物都被交付}给祂的意义，正如\BibleBook{玛窦}{福音}所载，祂继续说：\BibleQuote{凡劳苦和负重担的，你们都到我跟前来罢，我要使你们安息} \BibleRef{玛 11:28}。是的，你们\Quote{被交付}给我，是为使劳苦者得安息，使死者得生命。这与\BibleBook{若望}{福音}所记载的相符合：\BibleQuote{父爱子，并把一切交在祂手中} \BibleRef{若 3:35}。交付了，为的是，如同万物是藉着祂而造，同样万物也要在祂内得以更新。因为万物并非\Quote{交付}给祂，好让贫穷的祂变为富有；祂接受万物，也并非为了取得以前所缺的能力：这是决不可能的。而是为要以救主的身份，使万物恢复正常。因为这样才恰当：万物在起初是\Quote{藉着祂}而有，后来万物却应\Quote{在祂内}（注意用词的改变）得以复原 \BibleRef{若 1:3}。因为在起初，万物是\Quote{藉着}祂而受造；但后来万物都已堕落，圣言便成了血肉，穿上了人性，为的是\Quote{在祂内}使一切复原。祂自己受苦，赐我们安息；祂自己饥饿，却喂养我们；祂下降阴府，为将我们从那里领回。例如，在万物受造之时，创造是藉着一道成命完成的，如\Quote{[大地]生出}，\Quote{有光} \BibleRef{创 1:3}；但在复兴时，万物则适宜于\Quote{交付}给祂，为使祂成为人，使万物在祂内得以更新。因为人既在祂内，便获得生命：这正是圣言结合于人的原因，即令诅咒不再对人掌权。正因如此，我们在第七十一篇\BibleBook{}{圣咏}中读到为人类所作的祈求：\BibleQuote{天主，求你将审判权授予君王} \BibleRef{咏 71:1}：这祈求将悬在我们头上的死刑判决交付给圣子，使祂藉着为我们而死，在祂自己内为我们废除这判决。这正是祂在第八十七篇\BibleBook{}{圣咏}中自己所说的含义：\BibleQuote{你的忿怒重压着我} \BibleRef{咏 87:7}。因为祂承受了那本压在我们身上的忿怒，正如祂在第一百三十七篇\BibleBook{}{圣咏}中所说：\BibleQuote{上主，求你为我复仇} \BibleRef{咏 137:8}。\par

§3. \Emph{\Quote{万有}意指基督的救赎属性和能力。}\par

因此，我们可以理解万有都已交付给了救主，并且若有必要通过解释加深理解，那就是交付给祂的是祂先前所不拥有的。因为祂先前并非人，而是为了拯救人而成为人。圣言在起初并非是血肉，而是后来成了血肉\BibleRef{若 1:1}，在这血肉内，正如宗徒所说的，祂使那反对我们的仇恨得以和解\BibleRef{哥 1:20}，并废除了律法上诫命的规条，为将两者在祂内创造成一个新人，成就了和平，使双方在一个身体内与天主和好。然而，凡父所有的，也属于子，正如祂在 \BibleBook{若望}{福音}中所说的：\BibleQuote{凡父所有的一切，都是我的}\BibleRef{若 16:15}，这些表述是无可挑剔的。因为当祂成为祂原先所不是的时，\Quote{万有都交付}给了祂。但当祂愿意宣告祂与父的合一时，祂毫不保留地教导说：\BibleQuote{凡父所有的一切，都是我的。}人们不禁赞叹语言的精确。因为祂没有说\Quote{凡父所有的，祂都给了我}，以免显得祂曾一度不拥有这些；而是说\Quote{都是我的}。因为这些既在父的权能内，同样也在子的权能内。但我们必须转而探究\Quote{父所有的}是什么。如果是指创造，那么父在创造之前一无所有，并且似乎从创造中得到了某种附加的东西；但绝不可这样想。因为正如祂先于创造而存在，同样在创造之前祂就已拥有祂所拥有的，我们相信这些也都属于子\BibleRef{若 16:15}。因为如果子在父内，那么凡父所有的都属于子。因此，这一表述颠覆了异端者的悖谬，他们声称：\Quote{如果一切已交付给子，那么父就对所交付的东西不再有权柄，因为已立子代替自己。因为事实上，\BibleQuote{父不审判任何人，而将一切审判交给了子}？}\BibleRef{若 5:22}。但\BibleQuote{凡说恶言的，他们的口必被堵住}\BibleRef{咏 62:11}（因为虽然祂将一切审判交给了子，但祂并未因此被剥夺\OriginalTerm{主权}{lordship}；虽然说了万有由父交付给子，祂也并不减损其对万有的统御），他们显然将独生子与天主分离，而独生子按其本性是不可分离的，尽管他们疯狂地用言语分离祂；这些不敬虔的人不明白，光绝不可能离开太阳，因光本性就寓居于太阳内。因为我们不得不使用从可感知和熟悉的事物中抽取的贫乏比喻来表达我们的思想，因为贸然闯入不可理解的本性（天主的本性）是过于大胆的。\par

§4. \Emph{经文}\BibleRef{若 16:15}\Emph{，清楚地显示出圣子与圣父的本质关系。}\par

正如那照亮世界的来自太阳的光，心智健全的人绝不会认为它能离开太阳而如此运作，因为太阳的光与太阳本性相连；并且，如果光说：\Quote{我从太阳领受了照耀万物、并以我内的热力使它们生长强壮的能力}，没有人会疯狂到认为提及太阳是为了将光与其本性（即光）分离开；同样，虔敬要求我们领悟，圣言的\OriginalTerm{天主性体}{Divine Essence}按其本性与祂自己的父是合一的。因为我们面前的这段经文将把我们的问题置于最清晰的光照下，因为救主说过：\BibleQuote{凡父所有的一切，都是我的}；这表明祂永远与父同在。因为\Quote{凡祂所有的}表明父执掌\OriginalTerm{主权}{Lordship}，而\Quote{都是我的}则表明不可分离的合一。因此，我们必须领悟，在父内寓居着\OriginalTerm{永远性}{Everlastingness}、\OriginalTerm{永恒}{Eternity}、\OriginalTerm{不死}{Immortality}。这些在祂内并非作为偶性属性，而是如同在泉源中寓居于祂内，也寓居于圣子。因此，当你想领悟与子相关的事物时，就要了解父内所有的，因为这就是你必须相信存在于子内的。如果父是一件受造物或被造物，那么这些性质也属于子。如果可以说父\Quote{曾有一时不存在}或\Quote{从虚无中被造}，那么这些说法也可应用于子。但如果将这些属性归于父是不虔敬的，那么也当承认将它们归于子是不虔敬的。因为凡属于父的，也属于子。因为尊敬子的，就是尊敬派遣祂来的父；接受子的，就是接受父与祂同在，因为看见子的，就是看见了父\BibleRef{玛 10:40}。因此，父既不是受造物，子也不是；既然不能论及父说\Quote{曾有一时祂不存在}，也不能说\Quote{从虚无中被造}，那么同样论及子也是不合适的。相反，正如父的属性是\OriginalTerm{永远性}{Everlastingness}、\OriginalTerm{不死}{Immortality}、\OriginalTerm{永恒}{Eternity}，且非受造物，那么对于子我们也必须这样想。因为如经上所载\BibleRef{若 5:26}\BibleQuote{正如父在自己内自有生命，同样祂也赐给子在自己内自有生命。}但祂使用\Quote{赐给}一词以指明赐予的父。同样，正如生命在父内，也在子内，以此显示子是不可分离且永恒的。这就是为何祂精确地说\Quote{凡父所有的}，就是为了这样提及父，以免祂被视为父本身。因为祂没有说\Quote{我是父}，而是说\Quote{凡父所有的}。\par

§5. \Emph{同一经文进一步解释。}\par

\OriginalTerm{独生圣子}{Only-begotten Son}因为祂的独生圣子，你们\OriginalTerm{亚略派}{Arians}啊，可能被祂的\OriginalTerm{圣父}{Father}称为\Quote{父}，但并非按你们在错误中所理解的意思，而是（虽为生祂的\OriginalTerm{圣父}{Father}之子，却是）\Quote{永远的父亲}\BibleRef{依 9:6}。因为必须不留任何揣测的余地给你们。好，祂藉先知说：\BibleQuote{有一个婴孩为我们诞生了，有一个儿子赐给了我们；他肩上担负着王权，他的名称为神奇的谋士，大能的天主，永远的父亲，和平的君王}\BibleRef{依 9:6}。因此，天主的\OriginalTerm{独生圣子}{Only-begotten Son}同时是永远的父亲、大能的天主和统治者。并且清楚地显示，凡是\OriginalTerm{圣父}{Father}所有的一切，都是祂的；正如\OriginalTerm{圣父}{Father}赋予生命，圣子同样能随其意愿使人生活。因为祂说：\BibleQuote{死者要听见天主子的声音，且要生活}\BibleRef{若 5:25}，而\OriginalTerm{圣父}{Father}与圣子的意愿和渴望是合一的，因为他们的本性也是合一而不可分割的。而\OriginalTerm{亚略派}{Arians}徒然折磨自己，因为他们不明白我们救主的话：\Quote{父所有的一切，都是我的}。因为从这段话，可以立刻推翻\Person{撒伯里乌}的迷惑，也将暴露我们现代犹太人的愚妄。正因如此，\OriginalTerm{独生圣子}{Only-begotten Son}，既在自己内有生命，一如\OriginalTerm{圣父}{Father}所有，也唯独认识父是谁，即因为祂在父内，父在祂内。因为祂是祂的肖像，因此，因为祂是祂的肖像，凡属于\OriginalTerm{圣父}{Father}的一切，都在祂内。祂是精确的印玺，在自己内显示\OriginalTerm{圣父}{Father}；祂是生活的圣言和真理，德能，智慧，我们的圣化者和救赎者\BibleRef{格前 1:30}。因为\BibleQuote{我们生活、行动、存在，都在祂内}\BibleRef{宗 17:28}，并且\BibleQuote{除了父，没有一个知道子是谁；除了子，没有一个知道父是谁}\BibleRef{路 10:22}。\par

§6. \Emph{《三圣颂》被亚略派错误解释。其真正意义。}\par

这些不敬虔的人怎敢冒昧讲说愚昧的话，他们本不该如此，因为他们不过是人，连地上的事物如何描述都无法知晓？但我为什么说\Quote{地上的事物}呢？让他们告诉我们他们自己的本性，如果他们能发现如何探究他们自己的本性的话？他们实在是鲁莽而固执，竟不战栗地擅自论断那些天使也切望窥探的事\BibleRef{伯前 1:12}，天使的本性与地位远超他们。因为有什么比\OriginalTerm{革鲁宾}{Cherubim}或\OriginalTerm{色辣芬}{Seraphim}更亲近天主的呢？然而他们，甚至不能看见祂，也不能站立，甚至不用无遮盖的，而是仿佛以面纱遮脸，献上他们的赞颂，以不止息的唇舌只做一件事，就是藉着\OriginalTerm{三圣颂}{Trisagion}（\Quote{圣、圣、圣}）荣耀那神圣而不可言喻的本性。而且，那些受天主默示的众先知，就是特为这种神视而被拣选的人，从未向我们报告说：在颂念第一句\Quote{圣}时，声音高扬；在第二句时，则较低；而在第三句时，则极低——因此第一句表示主权，第二句表示从属，第三句标示更低下的等级。但让这些憎恨天主和愚昧之人的蠢话滚开吧。因为所赞颂、尊崇和朝拜的\OriginalTerm{天主圣三}{Triad}，乃是一体而不可分割的，且无等级之分（\OriginalTerm{无等级}{\GreekText{ἀσχηματιστός}}）。它结合而不相混淆，一如\OriginalTerm{一元}{Monad}也区分而不相分离。因为那些可敬的活物三次献上赞颂，说\BibleQuote{圣、圣、圣}，这证明了三个\OriginalTerm{位格}{Subsistences}是完美的，正如说\BibleQuote{主}时，他们宣认了同一本质\BibleRef{默 4:8}。所以，凡贬抑天主独生圣子的，就是亵渎天主，毁谤祂的完美并指控祂不完美，因而使自己担当最严厉的惩罚。因为亵渎任何一个\OriginalTerm{位格}{Subsistences}的，在今世和来世，都不得赦免。但天主能开启他们心灵的眼目，使其瞻仰\OriginalTerm{义德的太阳}{Sun of Righteousness}，好使他们认识了那先前所轻视的之后，能以坚定不移的虔诚之心，与我们一同光荣祂，因为国度属于祂，即属于圣父、圣子及圣神，从今直到永远。阿们。\par

\SourceRecord{Translated by Archibald Robertson.；Nicene and Post-Nicene Fathers, Second Series；Vol. 4.；Edited by Philip Schaff and Henry Wace.；Buffalo, NY: Christian Literature Publishing Co.,；1892.；<http://www.newadvent.org/fathers/2805.htm>.}
